\documentclass[11pt,a4paper]{article}
\usepackage[T1]{fontenc}
\usepackage{lmodern,amsmath,amssymb,amsthm,microtype,booktabs}
\usepackage[margin=25mm]{geometry}
\usepackage[colorlinks=true,urlcolor=blue,linkcolor=blue]{hyperref}
\newtheorem{proposition}{Proposition}
\newcommand{\dd}{\mathop{}\!\mathrm d}
\newcommand{\E}{\mathcal E}
\title{Joint endpoint control and central cancellation\\in the recoil Barnes integral}
\author{Alan Day \and Chengheng\thanks{AI research assistant (OpenAI).}}
\date{Working research note --- 30 October 2026}
\begin{document}\maketitle
\begin{abstract}
An optimized left endpoint and the exact Gamma modulus identities
control the negative vertical ray in a joint large-mass regime.
Reflection and Stirling estimates control the positive tail relative
to the sum of the two Bessel branch magnitudes. However, a central
fixed-height window has an exponentially larger absolute norm than
this envelope. Thus endpoint control does not complete the joint
asymptotic theorem: the remaining estimate must preserve phase.
These are self-reviewed analytic estimates and finite-point numerical
checks, not a finite-\(G\) holographic dictionary.
\end{abstract}

\section{Object, scale and remaining question}
Set the de Sitter radius to one, \(u=v=0\), \(\omega=i\alpha\),
\(K\ge1\), \(\alpha=Ky\), \(0\le y\le Y<\infty\), and
\(M=\nu^2/(2K)\). The fixed-mass \(\kappa\to\infty\) limit has already
been taken. This note concerns an internal chiral graph integral;
no physical external state, operator algebra, insertion normalization
or finite-\(G\) OTOC is assigned to it. Define
\[
 b=\tfrac12-i(M-\nu),\qquad c=\tfrac12-i(M+\nu),
\]
\begin{align}
 Q_\nu(t)&=e^{-i\pi t}
 \frac{\Gamma(t-i\nu/2)^4}{\Gamma(-i\nu/2)^4}
 \frac{\Gamma(b-t)\Gamma(c)}{\Gamma(b)\Gamma(c+t)}
 \Gamma(\tfrac12-t+\alpha/2)\Gamma(\tfrac12-t-\alpha/2),\\
 U(t)&=e^{-2\pi it}(K/2)^{2t}
 \Gamma(\tfrac12-t+\alpha/2)\Gamma(\tfrac12-t-\alpha/2).
\end{align}
The contour runs upward from \(-N-i\infty\) to \(-N+i\),
then right. Its integral includes \(1/(2\pi i)\). Previously,
the auxiliary integral was evaluated as
\[
 J=-K[e^{i\pi\alpha}K_\alpha(K)+i\pi I_\alpha(K)].
\]
Use the positive scale
\[
 \E=K[K_\alpha(K)+\pi I_\alpha(K)],\qquad
 F(y)=\sqrt{1+y^2}-y\operatorname{arsinh}y.
\]
For negative \(\alpha\), evenness of the complete expression permits
using \(|\alpha|\) in this envelope. No estimate below divides by \(J\).
Its zeros are genuine obstructions to that error convention.
Our hierarchy is \(K\to\infty\), \(\rho=K^3/\nu^2\le\rho_Y\),
with a sufficiently small fixed positive \(\rho_Y\).
Statements about convergence of corrections would additionally need
\(\rho\to0\); this note does not prove the full joint expansion.

\section{A multiplier estimate on growing compact regions}
Fix \(C<\infty\). On \(|t|\le CK\), the quotient obtained by canceling
the two small Gamma factors admits a continuous logarithm, with
\begin{equation}\label{eq:mult}
 \log(Q_\nu/U)=\frac{4i(t^2-t)}{\nu}+O_C(\rho).
\end{equation}
This is a growing-region estimate, not the direct use of a fixed-shift
Gamma asymptotic. Here and below constants may depend on \(Y,C\).

For completeness, apply sectorial Stirling to each large argument and
Taylor expand its shift, using
\[
 \log\Gamma(z+w)-\log\Gamma(z)
 =w\log z+\frac{w(w-1)}{2z}
 +O\!\left(\frac{|w|+|w|^2+|w|^3}{|z|^2}\right).
\]
The segments \(z+\theta w\) stay in a common closed sector away from
the negative real axis, since \(K/\nu\to0\).
The four numerator factors give \(4i(t^2-t)/\nu\).
The leading logarithms combine to
\(2t\log(K/2)-2\pi it\); the remaining \(b,c\) terms are bounded
by a constant times
\[
 \frac{K^2|t|+K|t|^2+|t|^3+K}{\nu^2}.
\]
This proves \eqref{eq:mult}, with the logarithm fixed continuously
from zero shift after the small factors have been canceled.
The first phase must be retained: \(K^2/\nu\) need not tend to zero.

At a real \(t=-N=O(K)\) this leading term is purely imaginary.
At \(t=x+i\) its real part is \(-4(2x-1)/\nu\).
Consequently the left real base and the positive base satisfy
\[
 |Q_\nu(-N)|\le e^{C_Y\rho}|U(-N)|,\qquad
 |Q_\nu(x+i)|\le e^{C_Y\rho}|U(x+i)|\quad (x\ge\tfrac12,\ x=O_Y(K)).
\]
On any fixed-height central strip, the modulus comparison is two-sided
with factors \(\exp[\pm C_Y(\rho+K/\nu)]\).

\section{The optimized left base and its vertical ray}
Choose
\[
 L=\frac K2\sqrt{1+y^2},\qquad N=\lceil L\rceil.
\]
For sufficiently large \(K\) depending only on \(Y\),
\(N>\alpha/2+1/2\), and both small Gamma real parts are bounded below
by \(c_YK>0\). Moving the left corner to this integer crosses no poles:
the small-Gamma pole families start to its right, the numerator poles
are at height \(\nu/2>1\), and the \(b-t\) poles have positive real part.
The lower closure vanishes by the fixed-parameter vertical decay.
Stirling on the positive real axis yields
\begin{equation}\label{eq:base}
 |U(-N)|=2\pi\exp[-KF(y)+O_Y(K^{-1})].
\end{equation}
Indeed, the leading exponent as a function of \(n\) is
\[
 g(n)=(n+\alpha/2)\log\frac{2(n+\alpha/2)}K
 +(n-\alpha/2)\log\frac{2(n-\alpha/2)}K-2n.
\]
At \(n=L\), \(g=-KF\), \(g'=0\), and \(g''=O_Y(K^{-1})\).
Thus rounding to \(N\) changes the exponent by only \(O_Y(K^{-1})\).

A lower bound for \(\E\) valid even when \(y\to0\) follows directly from
the standard positive integral \cite{Kint}:
\[
 K_\alpha(K)=\int_0^\infty e^{-K\cosh u}\cosh(\alpha u)\,\dd u.
\]
Put \(u_*=\operatorname{arsinh}y\).
On \([u_*,u_*+K^{-1/2}]\), the exponent
\(-\cosh u+yu\) is at least \(-F(y)-C_Y/K\).
Using \(\cosh(\alpha u)\ge e^{\alpha u}/2\) gives
\begin{equation}\label{eq:env}
 \E\ge K K_\alpha(K)\ge c_Y\sqrt K\,e^{-KF(y)},\qquad
 \frac{|Q_\nu(-N)|}{\E}\le C_YK^{-1/2}e^{C_Y\rho}.
\end{equation}

\begin{proposition}
Let \(B=M-\nu>0\), \(C_0=M+\nu\), and
\(\lambda=2\pi-N(B^{-1}+C_0^{-1})\).
For every \(s\ge0\), with \(N>\alpha/2+1/2\),
\begin{equation}\label{eq:vertical}
 |Q_\nu(-N-is)|\le |Q_\nu(-N)|e^{-\lambda s}.
\end{equation}
In the hierarchy above, \(\lambda\ge\pi\) eventually, uniformly in
\(y\in[0,Y]\), and
\[
 \frac1{\E}\int_{s_0}^\infty |Q_\nu(-N-is)|\,\dd s
 \le C_YK^{-1/2}e^{C_Y\rho-\pi s_0}.
\]
\end{proposition}

\begin{proof}
Gamma recurrence and the standard moduli \cite{Gamma} give
\begin{align*}
 |\Gamma(-N-iw)|^2
 &=\frac{\pi}{w\sinh(\pi w)\prod_{j=1}^N(j^2+w^2)},\quad w>0,\\
 |\Gamma(N+\tfrac12+iw)|^2
 &=\frac{\pi\prod_{j=0}^{N-1}[(j+\tfrac12)^2+w^2]}{\cosh(\pi w)}.
\end{align*}
The reciprocal negative half-integer formula follows by recurrence.
In the ratio to the base \(s=0\), the explicit exponential gives
\(e^{-\pi s}\). The four negative-integer Gamma factors give at most
\(e^{-2\pi s}\), since their rational factors decrease and
\(\sinh(a+b)\ge e^b\sinh a\) for \(a,b>0\).
The two large half-integer hyperbolic factors contribute at most
\(e^{\pi s}\), because \(\log\cosh\) has derivative bounded by one.

Each polynomial ratio is bounded using
\(\sqrt{a^2+(B-s)^2}\le(1+s/B)\sqrt{a^2+B^2}\), and likewise
with \(C_0+s\). Their combined contribution is at most
\[
 (1+s/B)^N(1+s/C_0)^N\le e^{Ns(B^{-1}+C_0^{-1})}.
\]
The remaining small Gamma ratios are at most one, by
\(|\Gamma(a+is)|\le\Gamma(a)\) for \(a>0\), a consequence of Euler's
integral and the triangle inequality. This proves
\eqref{eq:vertical}. Finally \(N(B^{-1}+C_0^{-1})=O_Y(K^2/\nu^2)\),
and integration gives the stated bound.
\end{proof}

The short connector from \(-N\) to \(-N+i\) has the same
\(O_Y(K^{-1/2})\) envelope bound, with factor
\(e^{C_Y(\rho+K/\nu)}\): use
\(|\Gamma(a-i\tau)|\le\Gamma(a)\), \(0\le\tau\le1\),
and \(|e^{-2\pi i(-N+i\tau)}|\le e^{2\pi}\).
These statements control the vertical pieces only. The negative
horizontal part of the central segment is not included.

\section{The positive base and the entire positive tail}
Set \(c_Y=8+2Y\), \(R=c_YK\). Reflect both small Gamma factors:
\[
 U(t)=\frac{\pi^2e^{-2\pi it}(K/2)^{2t}}
 {\cos\pi(t-\alpha/2)\cos\pi(t+\alpha/2)
 \Gamma(\tfrac12+t-\alpha/2)\Gamma(\tfrac12+t+\alpha/2)}.
\]
On \(t=x+i\), the cosines are bounded away from zero.
For \(x\in[R,R+1]\), let \(r=x/K\). Uniform Stirling gives
\[
 |U(x+i)|\le C_Ye^{-Kh_y(r)},\qquad
 h_y(r)=(r+y/2)\log[2(r+y/2)]
 +(r-y/2)\log[2(r-y/2)]-2r.
\]
The function is even and convex in \(y\) on \(|y|<2r\), hence
\(h_y(r)\ge h_0(r)=2r[\log(2r)-1]\).
Since \(h_0\) increases for \(r\ge c_Y\), the true supremum obeys
\[
 H_R:=\sup_{x\in[R,R+1]}|Q_\nu(x+i)|
 \le C_Ye^{C_Y\rho-Kh_0(c_Y)}.
\]
Use the previously established exact-recurrence tail lemma:
for \(D_Y=\max(2,Y^2/2)\) and
\(\nu^2\ge32(D_Y+1)K^3\),
\[
 \int_R^\infty|Q_\nu(x+i)|\,\dd x\le10(R+1)H_R.
\]
Together with \eqref{eq:env} and \(F(y)\le1\), it follows that
\begin{equation}\label{eq:positive}
 \frac1{\E}\int_R^\infty|Q_\nu(x+i)|\,\dd x
 \le C_Y\sqrt K\exp\{-[h_0(c_Y)-1]K+C_Y\rho\}.
\end{equation}
The positive endpoint is therefore exponentially small in this
normalization. Sampling \(H_R\) numerically is not used as proof
of its supremum bound.

\section{An exponential obstruction to central absolute bounds}
Fix \(y>1\). For \(t=Kx+i\), \(-y/2<x<y/2\), reflect just the
Gamma factor with negative real argument. Uniformly on compact
subintervals of this interval, positive-real Stirling then gives
\[
 |U(Kx+i)|\asymp_y e^{KA_y(x)},\qquad
 A_y(x)=u[\log(2u)-1]-v[\log(2v)-1],
 \quad u=y/2-x,\quad v=y/2+x.
\]
The periodic reflection prefactor is bounded above and below by
positive constants at height one; it cannot be discarded in the
complex integral. Direct differentiation gives
\[
 A_y'=-\log(y^2-4x^2),\qquad
 x_*=-\tfrac12\sqrt{y^2-1},\qquad A_y''(x_*)=-4\sqrt{y^2-1}.
\]
At this maximum,
\[
 A_y(x_*)=B(y):=y\operatorname{arcosh}y-\sqrt{y^2-1}.
\]
For any fixed small \(a>0\), the window
\(W_K=\{t=q+i:|q-Kx_*|\le a\sqrt K\}\) therefore satisfies
\[
 \int_{W_K}|U(t)|\,|\dd t|\ge c_{y,a}\sqrt K\,e^{KB(y)}.
\]
Let \(y_c=1.5088795615\ldots\) be the zero of \(F\).
For fixed \(y\ge y_c\), the standard large-order Bessel formulas
\cite{Bessel} give
\(\E\le C_y\sqrt K e^{-KF(y)}\).
Thus
\begin{equation}\label{eq:condition}
 \frac{\int_{W_K}|U|\,|\dd t|}{\E}
 \ge c_{y,a}e^{K\Delta(y)},\qquad \Delta(y)=B(y)+F(y)>0.
\end{equation}
Indeed \(\Delta'=\operatorname{arcosh}y-\operatorname{arsinh}y<0\)
and \(\Delta(y)\to0\) at infinity (its leading term is \(1/(2y)\)).
This proves positivity on \(y>1\), in particular on \(y\ge y_c\).
The two-sided multiplier estimate transfers \eqref{eq:condition}
to \(Q_\nu\), at cost \(e^{-C_Y(\rho+K/\nu)}\).

This does not refute a joint signed-amplitude asymptotic.
It shows that the fixed-height absolute norm can be exponentially
ill-conditioned relative even to the positive branch envelope.
Multiplying a pointwise polynomial remainder bound by this norm
does not yield the desired polynomial envelope error.
The same obstruction applies to the retained quadratic phase:
its modulus on this strip changes by only \(e^{O_Y(K/\nu)}\).

\section{Numerical checks and continuation}
All data use the exact inherited kernel and
\(\nu=32K^{3/2}\), so \(\rho=1/1024\) is fixed.
This tests bounds in a joint hierarchy, not convergence as
\(\rho\to0\).
For \(K=8,32,128,512\) and \(y=0,1,2\), we check the real-product
vertical identity against complex log-Gamma evaluation at ten
values of \(s\), including \(M-\nu,M,M+\nu,2M\).
We also check the reflection identity on three positive base
points per parameter pair. These are finite-point checks; the
proofs above provide the continuum estimates.

The central computation integrates absolute values only, over
\(|q+\sqrt3K/2|\le\sqrt K/4\), at height one.
Unit-or-shorter cells use 20/32 Gauss--Legendre nodes and 55/75
decimal digits; the last row is also checked at 95 digits.
The largest relative order change is below \(7.4\,10^{-27}\).
No signed full integral or cancellation-resolved relative error
is computed in this experiment.

\begin{center}
\begin{tabular}{rrrr}
\toprule
\(K\)&\(\int_{W_K}|U|/\E\)&\(\int_{W_K}|Q_\nu|/\E\)&
\(K^{-1}\log(\int_{W_K}|U|/\E)\)\\
\midrule
8&\(9.33072\,10^2\)&\(1.00797\,10^3\)&0.854810\\
16&\(5.65331\,10^3\)&\(5.94872\,10^3\)&0.540000\\
32&\(2.77261\,10^5\)&\(2.86796\,10^5\)&0.391647\\
64&\(7.95113\,10^8\)&\(8.12952\,10^8\)&0.320219\\
128&\(7.15645\,10^{15}\)&\(7.25825\,10^{15}\)&0.285209\\
\bottomrule
\end{tabular}
\end{center}
The analytic rate is \(\Delta(2)=0.2506620134229251\ldots\);
the finite-\(K\) rates include large constant-prefactor effects.
The 300 checks consist of 120 vertical identities, 120 vertical
inequalities, 12 decay-rate checks, 36 reflection identities,
10 order refinements and two precision refinements.
Base and local-multiplier samples are recorded separately,
without invented universal numerical constants.

The next task is a phase-preserving estimate of the central
integral, for example by a justified analytic branch decomposition
or contour deformation with explicit pole bookkeeping.
Only after a complete two-branch joint estimate should one restore
\(u,v\), transverse tails and the physical external interfaces.
Finite \(G\), state identification and a microscopic holographic
dual remain open. No priority claim is made.

\begin{thebibliography}{9}
\bibitem{Gamma} NIST DLMF,
\href{https://dlmf.nist.gov/5.4}{Sec.\ 5.4},
\href{https://dlmf.nist.gov/5.5}{Sec.\ 5.5},
\href{https://dlmf.nist.gov/5.11}{Sec.\ 5.11}:
Gamma moduli, recurrence, reflection and sectorial Stirling.
\bibitem{Kint} NIST DLMF,
\href{https://dlmf.nist.gov/10.32.E9}{Eq.\ 10.32.9}:
the positive real integral for \(K_\alpha\).
\bibitem{Bessel} NIST DLMF,
\href{https://dlmf.nist.gov/10.41}{Eqs.\ 10.41.3--4}:
large-order modified-Bessel expansions, used here only at fixed \(y>0\).
\end{thebibliography}
{\small\noindent Reproducible files:\par
\noindent\path{code/recoil_bessel_endpoints.py}\par
\noindent\path{results/stage4/recoil-bessel-endpoints/diagnostics.json}\par
\noindent Previous analytic input:\par
\noindent\path{notes/stage4-recoil-bessel-cancellation.md}\par}
\end{document}
